\documentclass{article}
\usepackage{amsmath, amssymb, amsthm}
\usepackage{hyperref}
\usepackage{geometry}
\geometry{margin=1in}

\title{Erd\H{o}s Problem \#985}
\date{Last updated: 2025-08-31}
\author{Source: erdosproblems.com}

\begin{document}
\maketitle

\noindent\textbf{Status:} OPEN \quad \textbf{No prize}

\noindent\textbf{Tags:} number theory

\medskip

\noindent\textbf{Problem Statement:}

Is it true that, for every prime $p$, there is a prime $q&#60;p$ which is a primitive root modulo $p$?




Artin conjectured that $2$ is a primitive root for infinitely many primes $p$, which Hooley \cite{Ho67b} proved assuming the Generalised Riemann Hypothesis. Heath-Brown \cite{He86b} proved that at least one of $2$, $3$, or $5$ is a primitive root for infinitely many primes $p$.




References

[He86b] Heath-Brown, D. R., Artin's conjecture for primitive roots . Quart. J. Math. Oxford Ser. (2) (1986), 27--38.

[Ho67b] Hooley, Christopher, On {A}rtin's conjecture . J. Reine Angew. Math. (1967), 209--220.

\noindent\textbf{Related OEIS sequences:} A002233, A219429, A103309, possible


\bigskip
\noindent\small{Source: \url{https://www.erdosproblems.com/985}}
\end{document}
